\section*{Anexo A. Descripción Detallada de la Cadena de Procesamiento para la Generación de Empalmes con Inteligencia Artificial}

\subsection*{A.1 Detección de Zonas de Inserción mediante Detección de Actividad Vocal}

La inserción del fragmento sintético debe realizarse en una región donde no exista
actividad vocal activa, de forma que la transición entre el contenido original y el
fragmento sintetizado ocurra en un punto de silencio natural. Con este propósito se
empleó Silero VAD \cite{silero_vad}, un modelo neuronal de detección de actividad
vocal (\textit{Voice Activity Detection}, VAD) que opera a nivel de \textit{frame} de
10~ms, ofreciendo mayor robustez ante variaciones en el nivel de ruido de fondo en
comparación con enfoques basados en umbrales de energía.

A partir del mapa de actividad vocal generado, se extrajeron las zonas de silencio
confirmado aplicando un margen de seguridad de 120~ms alrededor de cada segmento
de voz, con el fin de evitar cortes en los bordes del habla. Únicamente las zonas con
duración mínima de 400~ms fueron consideradas candidatas. La zona de empalme fue
seleccionada aleatoriamente entre las candidatas ubicadas en el rango central del audio,
comprendido entre el 25\% y el 75\% de la duración total, evitando que el modelo de
detección aprenda patrones asociados a posiciones fijas dentro de la señal.

\subsection*{A.2 Transcripción Automática del Contenido Hablado}

Con el objetivo de que el fragmento sintético reproduzca contenido lingüístico extraído
del propio audio original, se utilizó Whisper \cite{whisper} para transcribir dos
segmentos de voz diferenciados. El primer segmento, denominado zona~A, fue extraído
de la primera mitad del audio con una duración mínima de cinco segundos, y su
transcripción se utilizó para generar el fragmento que sería insertado como empalme.
El segundo segmento, zona~B, fue extraído de la segunda mitad del audio y su
transcripción se empleó para generar el segmento de referencia sintético independiente.

La separación entre los contenidos lingüísticos empleados en cada archivo de salida
responde a un criterio de diseño del conjunto de datos: al utilizar frases diferentes
provenientes de distintas regiones del audio, se busca que la única fuente de
discriminación disponible para el modelo sea la huella acústica introducida por el
proceso de síntesis, y no el contenido lingüístico del fragmento insertado.

\subsection*{A.3 Síntesis de Voz con Clonación del Locutor}

La generación de los fragmentos sintéticos se realizó mediante XTTS~v2
\cite{xtts_v2}, un modelo de conversión texto-voz (\textit{Text-to-Speech}, TTS) con
capacidad de clonación de voz a partir de una muestra de referencia. Como señal de
referencia se utilizaron diez segundos extraídos del segmento central del audio
original, región que generalmente presenta las condiciones acústicas más estables del
locutor. Aunque el modelo aproxima las características vocales del hablante, la señal
producida introduce patrones propios del proceso de síntesis que no están presentes en
la huella del dispositivo de adquisición original, constituyendo la fuente de información
que el sistema de detección busca explotar.

\subsection*{A.4 Construcción del Audio con Empalme Sintético}

El fragmento sintético correspondiente a la zona~A fue insertado en la zona de silencio
seleccionada, produciendo la estructura:

\begin{equation}
    \text{Audio}_{\text{fake}} = [\text{REAL}_A] + [\text{IA}] + [\text{REAL}_B]
    \label{eq:estructura}
\end{equation}

donde $\text{REAL}_A$ y $\text{REAL}_B$ corresponden a los segmentos del audio
original ubicados antes y después de la zona de inserción, e $\text{IA}$ representa el
fragmento sintetizado. Para suavizar las transiciones se aplicó un fundido cruzado
(\textit{crossfade}) de 80~ms en ambos puntos de empalme. Adicionalmente, se
realizó una normalización de amplitud basada en la raíz cuadrática media entre el
fragmento sintético y el contexto acústico del audio original:

\begin{equation}
    g = \frac{\text{RMS}(r[n])}{\text{RMS}(s[n])}
    \label{eq:rms_gain}
\end{equation}

donde $r[n]$ corresponde al contexto de referencia extraído del audio real y $s[n]$ al
fragmento sintetizado. El fragmento normalizado $s'[n] = g \cdot s[n]$ fue insertado en
la posición seleccionada con el fundido cruzado descrito, reduciendo la posibilidad de
que diferencias de amplitud triviales sean utilizadas como criterio de detección.

\subsection*{A.5 Adaptación a Condiciones de Grabación con Ruido de Fondo}

Para los audios adquiridos en condiciones ruidosas se incorporó una etapa adicional de
adaptación ambiental. Dado que XTTS~v2 produce señales en condiciones cercanas al
silencio, el fragmento sintético puede contrastar con el ruido de fondo del audio real,
constituyendo una fuente de evidencia trivial para la detección. El perfil de ruido
$\eta[n]$ fue extraído de las zonas de silencio detectadas por el sistema VAD y
mezclado con el fragmento sintético normalizado $s'[n]$ aplicando una relación
señal-ruido (SNR) configurable:

\begin{equation}
    s_{\text{noisy}}[n] = s'[n] + \alpha \cdot \eta[n]
    \label{eq:noisy}
\end{equation}

donde el factor de ganancia $\alpha$ se determina a partir de la SNR objetivo $\rho$
expresada en decibelios mediante:

\begin{equation}
    \alpha = \frac{\text{RMS}(s'[n])}{10^{\rho/20} \cdot \text{RMS}(\eta[n])}
    \label{eq:alpha}
\end{equation}

Este procedimiento permite que el fragmento insertado presente un perfil espectral más
coherente con el entorno acústico del dispositivo de grabación original, incrementando
la dificultad de la detección basada en características superficiales del ruido de fondo.

\subsection*{A.6 Archivos de Salida y Organización del Conjunto de Datos}

Por cada grabación procesada se generaron los siguientes archivos:

\begin{itemize}
    \item \textbf{Splicing\_IA.wav}: grabación completa con el fragmento sintético
    insertado, correspondiente a la clase de audio manipulado.

    \item \textbf{IA\_segment.wav}: segmento de referencia sintético generado a partir
    de la transcripción de la zona~B, con contenido lingüístico diferente al del empalme,
    destinado a caracterizar la huella acústica de la síntesis de forma aislada.

    \item \textbf{metadata.json}: parámetros del proceso de generación, incluyendo
    tiempos de inserción, transcripciones de las zonas A y B, regiones de actividad vocal
    detectadas por VAD y parámetros del empalme aplicado.
\end{itemize}

Los archivos fueron organizados por dispositivo de grabación, de forma que cada
subdirectorio contiene únicamente las manipulaciones derivadas de grabaciones
adquiridas con un mismo equipo. Esta organización facilita la realización de experimentos
que evalúen la capacidad del sistema para generalizar hacia dispositivos no observados
durante el entrenamiento.

% ── Referencias del anexo ─────────────────────────────────────────────────────
\begin{thebibliography}{9}

\bibitem{silero_vad}
Silero Team, ``Silero VAD: pre-trained enterprise-grade Voice Activity Detector (VAD),
Number Detector and Language Classifier,'' GitHub, 2021.
[Online]. Available: \url{https://github.com/snakers4/silero-vad}

\bibitem{whisper}
A.~Radford, J.~W.~Kim, T.~Xu, G.~Brockman, C.~McLeavey, and I.~Sutskever,
``Robust speech recognition via large-scale weak supervision,''
in \textit{Proc. 40th Int. Conf. Machine Learning (ICML)}, 2023.
[Online]. Available: \url{https://arxiv.org/abs/2212.04356}

\bibitem{xtts_v2}
E.~Casanova, J.~Weber, C.~Shulby, A.~C.~Junior, E.~Gölge, and M.~A.~Ponti,
``XTTS: a massively multilingual zero-shot text-to-speech model,''
in \textit{Proc. Interspeech}, 2024.
[Online]. Available: \url{https://arxiv.org/abs/2406.04904}

\end{thebibliography}
